\chapter{Graph Fundamentals}\label{ch:graphs}

A \defn{graph} is a set of \defn{vertices} (also called nodes) joined by \defn{edges}. Road maps, friendships, dependencies
between tasks, the cells of a maze, the states of a puzzle and the moves between them: an astonishing number of problems
become graph problems once you see them that way, and then a small toolbox of traversal algorithms answers them. This
chapter builds that toolbox. It introduces the vocabulary and the ways to store a graph, then the two fundamental searches,
depth-first and breadth-first, together with their direct applications: connected components, shortest paths in unweighted
graphs and grids, bipartiteness, cycles, and the special structure of graphs in which every vertex has exactly one outgoing
edge. It closes with two techniques for searching huge implicit graphs: bidirectional search and A*.

\begin{chapterintro}{\bronze\ \silver\ \gold\ \ioi}{Chapters~\ref{ch:complexity} and~\ref{ch:ds-basics}.}
  \item speak the language of graphs, and choose between adjacency lists, matrices and edge lists;
  \item run DFS and BFS, and know exactly what each one guarantees;
  \item find components, shortest paths in grids, escape routes and 2-colourings;
  \item detect and output cycles in undirected and directed graphs;
  \item analyse functional graphs, and find cycles with $\Oh{1}$ memory;
  \item search huge implicit graphs with bidirectional BFS and A*.
\end{chapterintro}

% ======================================================================================================================
\section[Vocabulary]{Vocabulary \hfill \bronze}\label{sec:graph-vocab}

We write $n$ for the number of vertices and $m$ for the number of edges. An edge of an \emph{undirected} graph joins two
vertices symmetrically; an edge of a \emph{directed} graph (a \defn[arc]{arc}) goes from a tail to a head. A graph is
\emph{weighted} when every edge carries a number, its length or cost (Figure~\ref{fig:graphkinds}). Unless a problem says
otherwise, assume a \emph{simple} graph: no \emph{self-loops} (edges from a vertex to itself) and no \emph{parallel edges}
(two edges joining the same pair). Contest inputs, however, often contain both, and the algorithms of this chapter handle
them with at most a line of care.

\begin{figure}[htbp]
\centering
\begin{tikzpicture}
  \begin{scope}
    \node[vtx] (a) at (0, 1.2) {0}; \node[vtx] (b) at (1.4, 1.9) {1}; \node[vtx] (c) at (2.6, 1) {2};
    \node[vtx] (d) at (1.2, 0) {3}; \node[vtx] (e) at (3.3, 2.2) {4}; \node[vtx, dim] (f) at (3.5, 0) {5};
    \draw[ed] (a) -- (b) -- (c) -- (d) -- (a) (b) -- (e);
    \node[lbl] at (1.8, -0.75) {undirected; vertex 5 is isolated};
  \end{scope}
  \begin{scope}[xshift=5.3cm]
    \node[vtx] (a) at (0, 1.2) {0}; \node[vtx] (b) at (1.4, 1.9) {1}; \node[vtx] (c) at (2.6, 1) {2};
    \node[vtx] (d) at (1.2, 0) {3};
    \draw[ded] (a) -- (b); \draw[ded] (b) -- (c); \draw[ded] (c) -- (d); \draw[ded] (d) -- (a); \draw[ded] (a) -- (c);
    \node[lbl] at (1.3, -0.75) {directed; cycle $0 \to 1 \to 2 \to 3 \to 0$};
  \end{scope}
  \begin{scope}[xshift=10.3cm]
    \node[vtx] (a) at (0, 1.2) {0}; \node[vtx] (b) at (1.6, 1.9) {1}; \node[vtx] (c) at (2.8, 0.8) {2};
    \node[vtx] (d) at (1.2, 0) {3};
    \draw[ed] (a) -- node[wl] {4} (b) (b) -- node[wl] {1} (c) (c) -- node[wl] {7} (d) (d) -- node[wl] {2} (a)
      (a) -- node[wl] {5} (c);
    \node[lbl] at (1.4, -0.75) {weighted};
  \end{scope}
\end{tikzpicture}
\caption{An undirected graph, a directed graph and a weighted graph.}
\label{fig:graphkinds}
\end{figure}

\begin{itemize}
  \item A \defn{path} is a sequence of vertices in which consecutive vertices are joined by edges; it is \emph{simple} if no
        vertex repeats. A \defn{cycle} is a path of at least one edge that returns to its start (in an undirected simple graph,
        of length at least 3).
  \item An undirected graph is \defn{connected} if every two vertices are joined by a path. Its maximal connected pieces are
        its \defn{connected components}; every vertex belongs to exactly one.
  \item The \defn{degree} of a vertex is the number of edges at it. Every edge adds one to the degrees of both its ends, so
        the degrees sum to $2m$ (the \emph{handshake lemma}); in particular, the number of vertices of odd degree is even.
        In a directed graph, each vertex has an in-degree and an out-degree, and both sum to $m$.
  \item A \defn{tree} is a connected graph without cycles. The following are equivalent for a graph with $n$ vertices: it is
        a tree; it is connected with $n - 1$ edges; it is acyclic with $n - 1$ edges; any two vertices are joined by exactly
        one simple path. A \defn{forest} is a graph whose components are trees.
  \item A \defn{DAG} is a directed acyclic graph. DAGs are exactly the graphs that have a topological order
        (Chapter~\ref{ch:dsu}), and they are the natural shape of every dynamic programming (Part~V).
  \item A graph is \defn{bipartite} if its vertices split into two sides with every edge crossing between them
        (Section~\ref{sec:bipartite}).
\end{itemize}

Many graphs in contests are \emph{implicit}: the vertices are the cells of a grid, the arrangements of a puzzle, the numbers
up to $10^6$, and the edges are the legal moves between them. Such graphs are never stored; the search generates the
neighbours of a vertex when it needs them.

% ======================================================================================================================
\section[Representing a graph]{Representing a graph \hfill \bronze}\label{sec:graph-rep}

Three representations cover almost every need (Figure~\ref{fig:graphrep}).

\begin{itemize}
  \item The \defn{adjacency list} stores, for every vertex, the list of its neighbours. It takes $\Oh{n + m}$ memory, and
        visiting the edges of a vertex costs its degree, so a whole traversal costs $\Oh{n + m}$. It is the default.
  \item The \defn{adjacency matrix} is an $n \times n$ table with a 1 (or the edge weight) where an edge exists. It answers
        ``is there an edge from $u$ to $v$?'' in $\Oh{1}$ but needs $\Oh{n^2}$ memory and $\Oh{n}$ to list a vertex's
        neighbours, so it suits only small or dense graphs: Floyd--Warshall, bitmask DP, $n$ up to a few thousand.
  \item The \defn{edge list} is just the array of edges. It is what Kruskal's algorithm sorts and what Bellman--Ford relaxes,
        and it is the natural format when edges are processed in some order rather than explored from vertices.
\end{itemize}

\begin{figure}[htbp]
\centering
\begin{tikzpicture}
  \begin{scope}
    \node[vtx] (a) at (0, 1.4) {0}; \node[vtx] (b) at (1.5, 1.9) {1}; \node[vtx] (c) at (2.4, 0.6) {2};
    \node[vtx] (d) at (1.0, -0.1) {3}; \node[vtx] (e) at (-0.4, -0.1) {4};
    \draw[ed] (a) -- (b) (a) -- (d) (b) -- (c) (b) -- (d) (d) -- (e);
    \node[lbl] at (1.0, -0.85) {the graph};
  \end{scope}
  \begin{scope}[xshift=4.0cm, yshift=1.9cm]
    \foreach \vi/\lst in {0/{1, 3}, 1/{0, 2, 3}, 2/{1}, 3/{0, 1, 4}, 4/{3}} {
      \node[font=\sffamily\small\bfseries, text=fgblue] at (0, -\vi*0.5) {\vi};
      \node[font=\ttfamily\small, anchor=west] at (0.35, -\vi*0.5) {\lst};
    }
    \node[lbl] at (0.9, -2.75) {adjacency list};
  \end{scope}
  \begin{scope}[xshift=7.3cm, yshift=1.9cm]
    \foreach \rw/\bits in {0/{0,1,0,1,0}, 1/{1,0,1,1,0}, 2/{0,1,0,0,0}, 3/{1,1,0,0,1}, 4/{0,0,0,1,0}} {
      \node[font=\sffamily\scriptsize\bfseries, text=fgblue] at (-0.45, -\rw*0.44) {\rw};
      \foreach \bv [count=\ci from 0] in \bits {
        \ifnum\bv=1 \node[cellc, cool, minimum width=4.4mm, minimum height=4.4mm, font=\ttfamily\scriptsize] at (\ci*0.44,
          -\rw*0.44) {1};
        \else \node[cellc, minimum width=4.4mm, minimum height=4.4mm, font=\ttfamily\scriptsize, text=fggray] at
          (\ci*0.44, -\rw*0.44) {0}; \fi
      }
    }
    \foreach \ci in {0,...,4} \node[font=\sffamily\scriptsize\bfseries, text=fgblue] at (\ci*0.44, 0.42) {\ci};
    \node[lbl] at (0.88, -2.75) {adjacency matrix};
  \end{scope}
  \begin{scope}[xshift=11.0cm, yshift=1.9cm]
    \foreach \ea/\eb [count=\kc from 0] in {0/1, 0/3, 1/2, 1/3, 3/4}
      \node[font=\ttfamily\small] at (0.6, -\kc*0.5) {\ea\ \eb};
    \node[lbl] at (0.6, -2.75) {edge list};
  \end{scope}
\end{tikzpicture}
\caption{One undirected graph in the three standard representations. The adjacency list has one entry per edge end ($2m$
in total), the matrix has $n^2$ entries whatever the number of edges, and the edge list has $m$ entries.}
\label{fig:graphrep}
\end{figure}

Reading a graph into adjacency lists is the first few lines of countless solutions. Most inputs number the vertices from 1;
either convert to 0-based indices while reading or allocate $n + 1$ lists and ignore list 0, but do it the same way
everywhere.

\begin{code}
int n, m;
cin >> n >> m;
vvi adj(n + 1);  // vertices 1..n, as in the input
rep(i, m) {
    int u, v;
    cin >> u >> v;
    adj[u].push_back(v);
    adj[v].push_back(u);  // leave this line out for a directed graph
}
\end{code}

\usetemplate{Read graph (n, m, m edges)}{Ctrl+Alt+T, Shift+G}{cpgraph}

\usetemplate{Read weighted graph (n, m, edges u v w)}{Ctrl+Alt+T, Shift+E}{cpgraphw}

\begin{trick}[Flat adjacency arrays]
A \cc{vector<vector<int>>} makes $n$ separate allocations and scatters the lists across memory. For graphs with millions of
edges, the \emph{compressed sparse row} format is faster and smaller: count the degrees, turn them into starting offsets with
a prefix sum, and fill one array of $2m$ neighbours, so that the neighbours of $u$ are \cc{nb[start[u]]} up to
\cc{nb[start[u + 1] - 1]}. It is counting sort (Section~\ref{sec:radix}) applied to the edge list. On a random graph with
$10^6$ vertices and $4 \cdot 10^6$ edges, building nested vectors took 0.97\,s against 0.08\,s for the flat array, and a
BFS took 0.10\,s against 0.07\,s: the many small allocations, not the traversal, are what hurts
(Section~\ref{sec:memory}).
\end{trick}

% ======================================================================================================================
\section[Depth-first search]{Depth-first search \hfill \silver}\label{sec:dfs}

\defn{Depth-first search} (DFS) explores a graph the way you would explore a maze with a ball of string: always go deeper
along an unexplored edge, and when there is none, back up to the previous vertex and try its next edge. Implemented
recursively it is a few lines, visits every vertex reachable from the start exactly once, and looks at every edge at most
twice, so it runs in $\Oh{n + m}$.

\begin{pseudo}{Depth-first search}
\begin{algorithmic}[1]
\Procedure{DFS}{$u$}
  \State mark $u$ as visited; \ $\mathit{tin}[u] \gets \mathit{timer}{+}{+}$ \Comment{discovery time}
  \ForAll{edges $u \to v$}
    \If{$v$ is not visited}
      \State $\mathit{parent}[v] \gets u$; \ \Call{DFS}{$v$} \Comment{a tree edge}
    \EndIf
  \EndFor
  \State $\mathit{tout}[u] \gets \mathit{timer}{+}{+}$ \Comment{finishing time}
\EndProcedure
\end{algorithmic}
\end{pseudo}

The edges along which DFS discovers new vertices form the \defn{DFS tree} (a forest if the graph is not connected). The
\emph{discovery} and \emph{finishing} times nest like parentheses: the interval $[\mathit{tin}[v], \mathit{tout}[v]]$ of a
descendant lies inside that of its ancestor, and the intervals of unrelated vertices are disjoint. Relative to the tree, every
edge of a directed graph falls into one of four classes (Figure~\ref{fig:dfsclasses}):
\begin{itemize}
  \item \textbf{tree edges} discover a new vertex;
  \item \textbf{back edges} go to an ancestor that is still open (on the recursion stack): each one closes a cycle;
  \item \textbf{forward edges} go to a descendant that is already finished;
  \item \textbf{cross edges} go to a vertex in another, already finished branch.
\end{itemize}
In an \emph{undirected} graph there are only tree edges and back edges: an edge to an already finished vertex in another
branch would have been explored from that vertex while it was open, as a tree edge. This single fact is the basis of bridges,
articulation points and cycle detection in undirected graphs (Chapter~\ref{ch:connectivity}), while the four classes drive
cycle detection, topological sorting and strongly connected components in directed graphs.

\begin{figure}[htbp]
\centering
\begin{tikzpicture}
  \node[vtx] (v0) at (0, 0) {0};
  \node[vtx] (v1) at (-1.7, -1.3) {1};
  \node[vtx] (v2) at (-1.7, -2.8) {2};
  \node[vtx] (v3) at (1.7, -1.3) {3};
  \node[vtx] (v4) at (1.7, -2.8) {4};
  \node[vtx] (v5) at (1.7, -4.3) {5};
  \draw[bded] (v0) -- (v1); \draw[bded] (v1) -- (v2); \draw[bded] (v0) -- (v3); \draw[bded] (v3) -- (v4);
  \draw[bded] (v4) -- (v5);
  \draw[->, line width=1.2pt, draw=fgverm, dashed] (v2) to[out=165, in=180, looseness=1.3] (v0);
  \draw[->, line width=1.2pt, draw=fgverm, dashed] (v5) to[out=15, in=0, looseness=1.3] (v3);
  \draw[->, line width=1.2pt, draw=fggreen, dash pattern=on 5pt off 2pt] (v0) to[bend left=12] (v2);
  \draw[->, line width=1.2pt, draw=fggray, dotted] (v3) -- (v2);
  \begin{scope}[xshift=4.6cm, yshift=-1.0cm]
    \draw[bded] (0, 0) -- (0.8, 0); \node[lbl, anchor=west] at (0.95, 0) {tree edge};
    \draw[->, line width=1.2pt, draw=fgverm, dashed] (0, -0.5) -- (0.8, -0.5);
    \node[lbl, anchor=west] at (0.95, -0.5) {back edge (to an open ancestor)};
    \draw[->, line width=1.2pt, draw=fggreen, dash pattern=on 5pt off 2pt] (0, -1.0) -- (0.8, -1.0);
    \node[lbl, anchor=west] at (0.95, -1.0) {forward edge (to a finished descendant)};
    \draw[->, line width=1.2pt, draw=fggray, dotted] (0, -1.5) -- (0.8, -1.5);
    \node[lbl, anchor=west] at (0.95, -1.5) {cross edge (to a finished other branch)};
    \node[anchor=north west, font=\sffamily\footnotesize, text=fgink!85, inner sep=0pt] at (0, -2.0) {%
      \begin{tabular}{@{}l c c c c c c@{}}
        vertex & 0 & 1 & 2 & 3 & 4 & 5 \\ \hline
        discovered & 1 & 2 & 3 & 6 & 7 & 8 \\
        finished & 12 & 5 & 4 & 11 & 10 & 9 \\
      \end{tabular}};
    \node[notel, anchor=north west, text width=5.8cm] at (0, -3.35) {adjacency lists $0\colon 1, 2, 3$; \ $1\colon 2$;
      \ $2\colon 0$; \ $3\colon 2, 4$; \ $4\colon 5$; \ $5\colon 3$};
  \end{scope}
\end{tikzpicture}
\caption{DFS from vertex 0 on a directed graph and the classification of its edges. The back edges $2 \to 0$ and $5 \to 3$
each close a cycle. The intervals from discovery to finishing time (table) nest exactly as the tree does: $[2, 5]$ and
$[6, 11]$ lie inside $[1, 12]$ and are disjoint from each other.}
\label{fig:dfsclasses}
\end{figure}

\begin{problem}{Building Roads (CSES)}
A country has $n$ cities and $m$ roads. Add the minimum number of new roads so that every city can reach every other, and
print them.
\solution Each connected component needs to be linked to the others, and $k$ components need exactly $k - 1$ new roads (each
road merges at most two components). Run a DFS from every city not yet visited: each start is a new component, and the start
vertices serve as representatives. Linking consecutive representatives gives the $k - 1$ roads.
\end{problem}

\begin{code}
void solve() {
    int n, m;
    cin >> n >> m;
    vvi adj(n + 1);
    rep(i, m) {
        int a, b;
        cin >> a >> b;
        adj[a].push_back(b);
        adj[b].push_back(a);
    }
    vector<bool> seen(n + 1, false);
    auto dfs = [&](auto&& self, int u) -> void {
        seen[u] = true;
        for (int v : adj[u])
            if (!seen[v]) self(self, v);
    };
    vi leader;  // one vertex from each connected component
    rep(v, 1, n + 1) if (!seen[v]) leader.push_back(v), dfs(dfs, v);
    print(sz(leader) - 1);  // a road between consecutive leaders
    rep(i, 1, sz(leader)) print(leader[i - 1], leader[i]);
}
\end{code}

\usetemplate{DFS (recursive, all components)}{Ctrl+Alt+T, D}{cpdfs}

\begin{pitfall}[Recursion depth]
On a path of $10^6$ vertices the recursion is $10^6$ calls deep, which needs tens of megabytes of stack. Codeforces and your
local builds (which use a 256\,MB stack) are fine, and CSES is fine too, but judges with the default 8\,MB or 1\,MB stack
crash with a segmentation fault that looks like a wrong answer. When in doubt, use an explicit stack (as in the flood fill of
Section~\ref{sec:grids}) or process vertices in BFS order, which suffices for most tree DPs.
\end{pitfall}

% ======================================================================================================================
\section[Breadth-first search]{Breadth-first search \hfill \silver}\label{sec:bfs}

\defn{Breadth-first search} (BFS) explores a graph in order of distance from the start: first the start, then all its
neighbours, then all their unvisited neighbours, and so on, layer by layer (Figure~\ref{fig:bfslayers}). A FIFO queue
produces this order. BFS runs in $\Oh{n + m}$ and, in an unweighted graph, finds the \emph{shortest} path from the start to
every vertex.

\emph{Why the distances are correct.} At every moment the queue holds vertices of at most two consecutive distances, $d$ and
$d + 1$, with all the $d$'s in front. (Initially it holds the start alone. Popping a vertex at distance $d$ appends
unvisited neighbours at distance $d + 1$ at the back, which keeps the property.) So vertices are popped in non-decreasing order
of their assigned distances, and a vertex is first discovered from a vertex of the smallest possible distance: by induction,
the distance assigned to each vertex is the length of a shortest path to it. Recording, for each vertex, the vertex it was
discovered from gives a \emph{BFS tree}, whose paths back to the start are shortest paths.

\begin{figure}[htbp]
\centering
\begin{tikzpicture}
  \foreach \rr/\cl in {0.62/fgverm, 1.52/fgorange, 2.42/fgblue} \draw[\cl, opacity=0.35, line width=0.8pt, dashed]
    (0, 0) circle (\rr);
  \node[vtx, hot] (s) at (0, 0) {0};
  \node[vtx] (a) at (1.05, 0.1) {1}; \node[vtx] (b) at (-0.6, 0.85) {1}; \node[vtx] (c) at (-0.5, -0.9) {1};
  \node[vtx] (d) at (1.9, 1.1) {2}; \node[vtx] (e) at (-1.6, 1.4) {2}; \node[vtx] (f) at (0.9, -1.7) {2};
  \node[vtx] (g) at (2.8, -0.3) {3}; \node[vtx] (h) at (-2.6, -1.2) {3};
  \draw[ed] (d) -- (g) (e) -- (h) (a) -- (f);
  \draw[bded] (s) -- (a); \draw[bded] (s) -- (b); \draw[bded] (s) -- (c); \draw[bded] (a) -- (d); \draw[bded] (b) -- (e);
  \draw[bded] (c) -- (f); \draw[bded] (f) -- (g); \draw[bded] (c) -- (h);
  \node[notel, anchor=west, text width=5.4cm] at (3.6, 0.3) {Vertices are labelled with their distance from the start.
    Blue arrows form a BFS tree: every vertex is reached from a vertex one layer closer. Grey edges join vertices of
    equal or adjacent layers and are never used for discovery.};
\end{tikzpicture}
\caption{BFS explores layer by layer. The queue holds at most two consecutive layers at a time, so the first discovery of a
vertex is along a shortest path.}
\label{fig:bfslayers}
\end{figure}

\usetemplate{BFS (shortest path, unweighted)}{Ctrl+Alt+T, B}{cpbfs}

\begin{problem}{Labyrinth (CSES)}
In an $n \times m$ labyrinth ($n, m \le 1000$) with walls, find a shortest path from \texttt{A} to \texttt{B} and print it as
a string of moves \texttt{L}, \texttt{R}, \texttt{U}, \texttt{D}, or report that none exists.
\solution BFS from \texttt{A} over the free cells. Instead of a parent cell, store the direction of the move that first
reached each cell; walking these directions backwards from \texttt{B} (undoing each move) leads to \texttt{A}, and reversing
the collected letters gives the path. The grid is never turned into an adjacency list: the neighbours of a cell are generated
from direction arrays.
\end{problem}

\begin{code}
void solve() {
    int n, m;
    cin >> n >> m;
    vector<string> g(n);
    cin >> g;
    const int dr[4] = {-1, 1, 0, 0}, dc[4] = {0, 0, -1, 1};
    const string name = "UDLR";
    vvi from(n, vi(m, -1));  // the direction that first reached each cell
    queue<pii> q;
    rep(r, n) rep(c, m) if (g[r][c] == 'A') q.push({r, c}), from[r][c] = 4;
    while (!q.empty()) {
        auto [r, c] = q.front();
        q.pop();
        rep(d, 4) {
            int nr = r + dr[d], nc = c + dc[d];
            if (nr < 0 || nr >= n || nc < 0 || nc >= m) continue;
            if (g[nr][nc] == '#' || from[nr][nc] != -1) continue;
            from[nr][nc] = d;
            q.push({nr, nc});
        }
    }
    int r = 0, c = 0;
    rep(i, n) rep(j, m) if (g[i][j] == 'B') r = i, c = j;
    if (from[r][c] == -1) {
        print("NO");
        return;
    }
    string path;  // walk back from B to A along the recorded directions
    while (from[r][c] != 4) {
        int d = from[r][c];
        path += name[d];
        r -= dr[d], c -= dc[d];
    }
    reverse(all(path));
    print("YES");
    print(sz(path));
    print(path);
}
\end{code}

\usetemplate{Grid BFS (shortest path on a grid)}{Ctrl+Alt+T, G}{cpgridbfs}

\subsection{Several sources at once}

To find, for every vertex, the distance to the \emph{nearest} of several sources (the nearest hospital, the first moment fire
reaches a cell), there is no need for one BFS per source. Put all the sources into the queue with distance 0 before the
search starts: the result is the same as a BFS from an imaginary super-source joined to all of them, which is exactly the
distance to the nearest source. This \defn{multi-source BFS} costs one ordinary BFS.

\begin{problem}{Monsters (CSES)}
You and several monsters are in a labyrinth. Each step, you move to an adjacent free cell, and so does every monster, in the
worst possible way for you. Escape through any border cell, never entering a cell that a monster can reach at the same time
or earlier.
\solution One multi-source BFS from all monsters computes the earliest time a monster can occupy each cell. A second BFS from
your start then enters a cell only if you arrive strictly earlier than every monster (Figure~\ref{fig:monsters}). Any border
cell that the second BFS reaches is an escape; its path is recovered from stored directions as in Labyrinth. The key insight
is that the monsters need not be simulated step by step: if you can reach a cell strictly before any monster, no monster
strategy can catch you there.
\end{problem}

\begin{figure}[htbp]
\centering
\begin{tikzpicture}[x=0.62cm, y=0.62cm]
  \foreach \rw/\cl in {0/{1,1,0,1,1,1,1,1}, 1/{1,0,0,0,0,0,0,1}, 2/{1,0,1,1,0,1,0,1}, 3/{1,0,0,0,0,1,0,1},
      4/{1,0,1,0,0,1,0,1}, 5/{1,1,1,0,1,1,0,1}, 6/{1,0,0,0,0,0,0,1}, 7/{1,1,1,1,1,0,1,1}} {
    \foreach \w [count=\ci from 0] in \cl {
      \ifnum\w=1 \fill[fgink!75] (\ci, -\rw) rectangle ++(1, -1);
      \else \draw[fgline!70, line width=0.3pt, fill=white] (\ci, -\rw) rectangle ++(1, -1); \fi
    }
  }
  \foreach \rw/\ci in {1/5, 1/6, 2/6, 3/6, 4/6, 5/6, 6/1, 6/2, 6/3, 6/4, 6/5, 6/6, 7/5}
    \fill[fggreen!25] (\ci + 0.04, -\rw - 0.04) rectangle ++(0.92, -0.92);
  \draw[line width=4pt, draw=fggreen!55, opacity=0.8, line join=round] (6.5, -3.5) -- (6.5, -6.5) -- (5.5, -6.5)
    -- (5.5, -8.0);
  \foreach \rw/\ci/\tv in {0/2/2, 1/2/1, 1/3/2, 1/4/3, 1/5/4, 1/6/5, 2/1/1, 2/4/4, 2/6/6, 3/1/2, 3/2/3, 3/3/4, 3/4/5,
      4/1/3, 4/3/5, 4/4/6, 4/6/8, 5/3/6, 5/6/9, 6/1/9, 6/2/8, 6/3/7, 6/4/8, 6/5/9, 6/6/10, 7/5/10}
    \node[font=\sffamily\scriptsize, text=fgverm!80!black] at (\ci + 0.5, -\rw - 0.5) {\tv};
  \node[font=\sffamily\small\bfseries, text=fgverm] at (1.5, -1.5) {M};
  \node[font=\sffamily\small\bfseries, text=fggreen!50!black] at (6.5, -3.5) {A};
  \draw[line width=1.2pt, draw=fgverm] (2.2, 0.35) -- (2.8, 0.95) (2.2, 0.95) -- (2.8, 0.35);
  \node[font=\sffamily\scriptsize, text=fgverm, anchor=west] at (2.95, 0.65) {exit lost};
  \node[font=\sffamily\scriptsize, text=fggreen!45!black, anchor=west] at (6.1, -8.35) {escape};
  \node[notel, anchor=west, text width=5.6cm] at (8.6, -2.6) {Numbers: the earliest time a monster can be in each cell
    (multi-source BFS from M). Green: the cells you reach strictly earlier. The top exit is reached by a monster at time 2,
    long before you could get there; the bottom exit is safe, reached at time 5 against the monster's 10.};
\end{tikzpicture}
\caption{Monsters as two BFS runs. You may enter a cell only if you get there strictly before any monster; the escape path
(thick green line) goes through such cells to the bottom border.}
\label{fig:monsters}
\end{figure}

\begin{code}
void solve() {
    int n, m;
    cin >> n >> m;
    vector<string> g(n);
    cin >> g;
    const int dr[4] = {-1, 1, 0, 0}, dc[4] = {0, 0, -1, 1};
    const string name = "UDLR";
    const int UNSEEN = 1e9;  // larger than any distance
    // multi-source BFS from all monsters: when can a monster reach a cell?
    vvi mon(n, vi(m, UNSEEN));
    queue<pii> q;
    rep(r, n) rep(c, m) if (g[r][c] == 'M') mon[r][c] = 0, q.push({r, c});
    auto bfs = [&](vvi& dist, vvi* from) {
        while (!q.empty()) {
            auto [r, c] = q.front();
            q.pop();
            rep(d, 4) {
                int nr = r + dr[d], nc = c + dc[d];
                if (nr < 0 || nr >= n || nc < 0 || nc >= m) continue;
                if (g[nr][nc] == '#' || dist[nr][nc] != UNSEEN) continue;
                int t = dist[r][c] + 1;
                if (from && t >= mon[nr][nc]) continue;  // monster is first
                dist[nr][nc] = t;
                if (from) (*from)[nr][nc] = d;
                q.push({nr, nc});
            }
        }
    };
    bfs(mon, nullptr);
    vvi me(n, vi(m, UNSEEN)), from(n, vi(m, -1));
    rep(r, n) rep(c, m) if (g[r][c] == 'A') me[r][c] = 0, q.push({r, c});
    bfs(me, &from);
    rep(r, n) rep(c, m) {
        bool border = r == 0 || r == n - 1 || c == 0 || c == m - 1;
        if (!border || me[r][c] == UNSEEN) continue;
        string path;  // reachable border cell: walk back to the start
        for (int y = r, x = c; me[y][x] > 0; ) {
            int d = from[y][x];
            path += name[d];
            y -= dr[d], x -= dc[d];
        }
        reverse(all(path));
        print("YES");
        print(sz(path));
        print(path);
        return;
    }
    print("NO");
}
\end{code}

\begin{trick}[BFS with 0--1 weights, and BFS on states]
If every edge has weight 0 or 1, a \emph{deque} replaces the queue: push the endpoint of a 0-edge to the front and of a
1-edge to the back. This \defn{0--1 BFS} keeps the two-layers invariant and runs in $\Oh{n + m}$ (Chapter~\ref{ch:shortest}).
And when the rules of movement depend on more than the position (keys collected, direction faced, remaining jumps), BFS runs
unchanged on the graph whose vertices are \emph{states} (position, extra information); the only question is whether the
number of states fits in time and memory.
\end{trick}

% ======================================================================================================================
\section[Grids and flood fill]{Grids and flood fill \hfill \silver}\label{sec:grids}

A grid is a graph whose vertices are cells and whose edges join neighbouring cells, usually the four orthogonal neighbours.
\defn{Flood fill} labels each connected region: start a DFS or BFS from any unlabelled free cell and label everything it
reaches (Figure~\ref{fig:floodfill}). Every cell is visited once, so labelling the whole grid takes $\Oh{nm}$.

\begin{figure}[htbp]
\centering
\begin{tikzpicture}[x=0.55cm, y=0.55cm]
  \foreach \rw/\ci/\rmid in {0/0/0, 0/1/0, 0/3/1, 0/4/1, 0/5/1, 0/6/1, 0/8/2, 0/9/2, 1/0/0, 1/1/0, 1/3/1, 1/6/1, 1/8/2,
      1/9/2, 2/3/1, 2/5/1, 2/6/1, 3/0/3, 3/1/3, 3/2/3, 3/5/1, 3/8/1, 3/9/1, 4/0/3, 4/2/3, 4/3/3, 4/5/1, 4/6/1, 4/7/1,
      4/9/1, 5/0/3, 5/1/3, 5/3/3, 5/5/1, 5/7/1, 5/8/1, 5/9/1} {
    \ifcase\rmid \def\fc{fgsky!45}\or\def\fc{fgorange!40}\or\def\fc{fgpink!40}\or\def\fc{fggreen!35}\fi
    \fill[\fc] (\ci, -\rw) rectangle ++(1, -1);
  }
  \foreach \rw/\ci in {0/2, 0/7, 1/2, 1/4, 1/5, 1/7, 2/0, 2/1, 2/2, 2/4, 2/7, 2/8, 2/9, 3/3, 3/4, 3/6, 3/7, 4/1, 4/4,
      4/8, 5/2, 5/4, 5/6}
    \fill[fgink!75] (\ci, -\rw) rectangle ++(1, -1);
  \draw[fgline!70, line width=0.3pt, step=1] (0, 0) grid (10, -6);
  \node[notel, anchor=west, text width=5.2cm] at (10.6, -3) {Four rooms, each found by one search started at its first
    free cell in reading order. The big orange room winds through gaps in the walls; the search follows it regardless of
    its shape.};
\end{tikzpicture}
\caption{Flood fill on a grid: dark cells are walls, and each colour is one connected region of free cells.}
\label{fig:floodfill}
\end{figure}

\begin{problem}{Counting Rooms (CSES)}
Count the rooms (connected regions of floor cells) in an $n \times m$ map, $n, m \le 1000$.
\solution Scan the cells; every unvisited floor cell starts a new room, which is then flooded. The code floods with an
explicit stack, marking each cell when it is pushed, so that no cell enters the stack twice and the recursion depth is no
concern, even for a $1000 \times 1000$ spiral.
\end{problem}

\begin{code}
void solve() {
    int n, m;
    cin >> n >> m;
    vector<string> g(n);
    cin >> g;
    const int dr[4] = {-1, 1, 0, 0}, dc[4] = {0, 0, -1, 1};
    int rooms = 0;
    rep(sr, n) rep(sc, m) {
        if (g[sr][sc] != '.') continue;
        rooms++;  // a new room: flood it with an explicit stack
        vector<pii> st{{sr, sc}};
        g[sr][sc] = '#';  // mark cells when they are pushed
        while (!st.empty()) {
            auto [r, c] = st.back();
            st.pop_back();
            rep(d, 4) {
                int nr = r + dr[d], nc = c + dc[d];
                if (nr < 0 || nr >= n || nc < 0 || nc >= m) continue;
                if (g[nr][nc] != '.') continue;
                g[nr][nc] = '#';
                st.push_back({nr, nc});
            }
        }
    }
    print(rooms);
}
\end{code}

\usetemplate{Flood fill (grid regions)}{Ctrl+Alt+T, F}{cpflood}

\begin{tip}[Grid habits]
Write the bounds check once, as the first line of the neighbour loop. For eight directions, use
\begin{code}
const int dr[8] = {-1, -1, -1, 0, 0, 1, 1, 1};
const int dc[8] = {-1, 0, 1, -1, 1, -1, 0, 1};
\end{code}
Encode a cell as the single number
$r \cdot m + c$ when a data structure (a DSU, a visited array, a hash set) wants one index. And remember that the
\emph{perimeter} of a region is the number of (cell, direction) pairs whose neighbour is outside the region, which a flood
fill can count as it goes: that is the whole of Icy Perimeter (USACO 2019 January Silver).
\end{tip}

% ======================================================================================================================
\section[Bipartite graphs]{Bipartite graphs and 2-colouring \hfill \silver}\label{sec:bipartite}

A graph is \defn{bipartite} if its vertices can be coloured with two colours so that every edge joins two different colours.
The characterisation is clean: \emph{a graph is bipartite if and only if it has no cycle of odd length}. One direction is
immediate: going around a cycle, the colours must alternate, so the cycle has even length. For the other, run BFS in every
component and colour each vertex by the parity of its distance. An edge between two vertices of the same colour joins two
vertices whose distances have the same parity (the distances of adjacent vertices differ by at most one, so they are equal);
together with the two BFS paths back to their lowest common ancestor in the BFS tree, that edge closes a cycle of odd length
(Figure~\ref{fig:bipartite}). So either the parity colouring works, or the graph has an odd cycle.

\begin{figure}[htbp]
\centering
\begin{tikzpicture}
  \begin{scope}
    \node[vtx, cool] (a) at (0, 0) {1}; \node[vtx, warn] (b) at (1.3, 0.8) {2}; \node[vtx, warn] (c) at (1.3, -0.8) {3};
    \node[vtx, cool] (d) at (2.6, 0.8) {4}; \node[vtx, cool] (e) at (2.6, -0.8) {5}; \node[vtx, warn] (f) at (3.9, 0) {6};
    \draw[ed] (a) -- (b) (a) -- (c) (b) -- (d) (c) -- (e) (b) -- (e) (d) -- (f) (e) -- (f);
    \node[lbl, text width=4.6cm, align=center] at (1.95, -1.65) {(a) BFS layers alternate colours;\\every edge crosses: bipartite};
  \end{scope}
  \begin{scope}[xshift=7.2cm]
    \foreach \kk/\cl in {0/cool, 1/warn, 2/cool, 3/warn, 4/cool} {
      \pgfmathsetmacro{\ang}{90 + \kk*72}
      \node[vtx, \cl] (p\kk) at (\ang:1.15) {\the\numexpr\kk+1\relax};
    }
    \draw[ed] (p0) -- (p1) (p1) -- (p2) (p2) -- (p3) (p3) -- (p4);
    \draw[hed] (p4) -- (p0);
    \node[lbl, text width=4.6cm, align=center] at (0, -1.95) {(b) an odd cycle: the red edge joins two\\vertices of the same
      colour};
  \end{scope}
\end{tikzpicture}
\caption{Two-colouring by BFS parity. In (a) the colouring succeeds. In (b), a 5-cycle, some edge must join equal colours:
graphs with an odd cycle are exactly the graphs that are not bipartite.}
\label{fig:bipartite}
\end{figure}

\begin{problem}{Building Teams (CSES)}
Divide $n$ pupils into two teams so that no two friends are in the same team, or report that it is impossible.
\solution BFS from every uncoloured pupil, giving each newly discovered pupil the colour opposite to the one it was discovered
from. If an edge ever joins two pupils of the same colour, there is an odd cycle and no division exists.
\end{problem}

\begin{code}
void solve() {
    int n, m;
    cin >> n >> m;
    vvi adj(n + 1);
    rep(i, m) {
        int a, b;
        cin >> a >> b;
        adj[a].push_back(b);
        adj[b].push_back(a);
    }
    vi team(n + 1, 0);  // 0 = not coloured yet, otherwise 1 or 2
    rep(s, 1, n + 1) {
        if (team[s]) continue;
        team[s] = 1;
        queue<int> q;
        q.push(s);
        while (!q.empty()) {
            int u = q.front();
            q.pop();
            for (int v : adj[u]) {
                if (!team[v]) team[v] = 3 - team[u], q.push(v);
                else if (team[v] == team[u]) {  // an odd cycle
                    print("IMPOSSIBLE");
                    return;
                }
            }
        }
    }
    print(vi(team.begin() + 1, team.end()));
}
\end{code}

\usetemplate{Bipartite check (2-coloring)}{Ctrl+Alt+T, 2}{cpbipartite}

\begin{recognize}
``Split into two groups such that each pair of rivals is separated'', ``two teams'', ``can the map be 2-coloured'', and
checkerboard arguments on grids (a grid graph is always bipartite: colour by $(r + c) \bmod 2$). Constraints of the form
``$x$ and $y$ in the \emph{same} group'' mixed with ``different groups'' are handled by a DSU with parity
(Chapter~\ref{ch:dsu}), and arbitrary boolean constraints on pairs by 2-SAT (Chapter~\ref{ch:connectivity}). Bipartite graphs
are also where matching theory lives (Chapter~\ref{ch:flows}).
\end{recognize}

% ======================================================================================================================
\section[Finding cycles]{Finding cycles \hfill \gold}\label{sec:cycles}

\paragraph{Undirected graphs.} In a DFS of an undirected graph, every non-tree edge goes to an ancestor
(Section~\ref{sec:dfs}), so a cycle exists exactly when some vertex $u$ has an edge to an already visited vertex $v$ other
than its parent. The cycle is $v$, the tree path down to $u$, and the edge back. Two details matter in practice. Parallel
edges to the parent must not count (they would make a 2-cycle), which is why the code skips the parent vertex rather than
the parent edge. And an edge from $u$ to a \emph{child} that was already explored can reappear through a parallel edge;
comparing depths distinguishes it from a real back edge, which goes \emph{up}.

\begin{problem}{Round Trip (CSES)}
Find a round trip, a route that visits at least three distinct cities and returns to its start without repeating a city, or
report that none exists.
\end{problem}

\begin{code}
void solve() {
    int n, m;
    cin >> n >> m;
    vvi adj(n + 1);
    rep(i, m) {
        int a, b;
        cin >> a >> b;
        adj[a].push_back(b);
        adj[b].push_back(a);
    }
    vi par(n + 1, 0), depth(n + 1, -1);  // depth -1: not visited yet
    vi cycle;
    // returns true once a cycle has been found and stored in cycle
    auto dfs = [&](auto&& self, int u) -> bool {
        for (int v : adj[u]) {
            if (v == par[u]) continue;  // back along the tree edge
            if (depth[v] == -1) {
                par[v] = u, depth[v] = depth[u] + 1;
                if (self(self, v)) return true;
            } else if (depth[v] < depth[u]) {  // an edge up to an ancestor
                for (int x = u; x != v; x = par[x]) cycle.push_back(x);
                cycle.push_back(v);
                cycle.push_back(u);
                return true;
            }  // otherwise a child reached again through a parallel edge
        }
        return false;
    };
    rep(s, 1, n + 1) {
        if (depth[s] != -1) continue;
        depth[s] = 0;
        if (!dfs(dfs, s)) continue;
        print(sz(cycle));
        print(cycle);
        return;
    }
    print("IMPOSSIBLE");
}
\end{code}

When only the \emph{existence} of a cycle matters, a DSU is even simpler: add the edges one by one, and the first edge whose
endpoints are already connected closes a cycle (Chapter~\ref{ch:dsu}).

\paragraph{Directed graphs.} In a directed graph, an edge to a visited vertex does not imply a cycle: cross and forward edges
reach finished vertices without closing anything. What closes a cycle is an edge to a vertex that is still \emph{open}, on the
current recursion path. So DFS keeps three states: white (not visited), grey (open, on the stack) and black (finished). An
edge into a grey vertex is a back edge, and the cycle is the stack segment from that vertex down to the current one
(Figure~\ref{fig:dircycle}). Conversely, if DFS finds no back edge, the finishing order is a topological order, so the graph is
acyclic: a directed graph has a cycle if and only if DFS finds a back edge.

\begin{figure}[htbp]
\centering
\begin{tikzpicture}[wht/.style={fill=white, draw=fgink}, gry/.style={fill=fgline!45, draw=fgink},
    blk/.style={fill=fgink!85, draw=fgink}]
  \node[vtx, fill=fgline!45, draw=fgink] (a) at (0, 0) {1};
  \node[vtx, fill=fgline!45, draw=fgink] (b) at (1.6, 0) {2};
  \node[vtx, fill=fgline!45, draw=fgink] (c) at (3.2, 0) {3};
  \node[vtx, fill=fgline!45, draw=fgink] (d) at (4.8, 0) {4};
  \node[vtx, fill=fgink!85, text=white, draw=fgink] (e) at (1.6, -1.4) {5};
  \node[vtx, fill=white] (f) at (4.8, -1.4) {6};
  \draw[bded] (a) -- (b); \draw[bded] (b) -- (c); \draw[bded] (c) -- (d);
  \draw[ded] (b) -- (e); \draw[ded] (c) -- (e);
  \draw[hded] (d) to[bend right=38] (b);
  \draw[ded, draw=fgline] (d) -- (f);
  \node[font=\sffamily\scriptsize, text=fgverm, above] at (3.2, 0.95) {edge into a grey vertex: cycle $2 \to 3 \to 4 \to 2$};
  \legendat{(6.4, 0.3)}{wht/white: not visited, gry/grey: open (on the stack), blk/black: finished}
\end{tikzpicture}
\caption{Three colours in a directed DFS. The search is at vertex 4 with the path $1 \to 2 \to 3 \to 4$ open. The edge
$3 \to 5$ led to a finished vertex (no cycle there), but the edge $4 \to 2$ enters an open vertex and closes a cycle.}
\label{fig:dircycle}
\end{figure}

\begin{problem}{Round Trip II (CSES)}
Find a directed cycle in a graph of $n \le 10^5$ cities and $m \le 2\cdot10^5$ one-way flights, and print it.
\end{problem}

\begin{code}
void solve() {
    int n, m;
    cin >> n >> m;
    vvi adj(n + 1);
    rep(i, m) {
        int a, b;
        cin >> a >> b;
        adj[a].push_back(b);
    }
    vi color(n + 1, 0), par(n + 1, -1);  // 0 white, 1 grey (on stack), 2 black
    vi cycle;
    auto dfs = [&](auto&& self, int u) -> bool {
        color[u] = 1;
        for (int v : adj[u]) {
            if (color[v] == 1) {  // an edge back into the current path
                cycle.push_back(v);
                for (int x = u; x != v; x = par[x]) cycle.push_back(x);
                cycle.push_back(v);
                reverse(all(cycle));
                return true;
            }
            if (color[v] == 0) {
                par[v] = u;
                if (self(self, v)) return true;
            }
        }
        color[u] = 2;  // finished: no cycle goes through u any more
        return false;
    };
    rep(s, 1, n + 1) {
        if (color[s] || !dfs(dfs, s)) continue;
        print(sz(cycle));
        print(cycle);
        return;
    }
    print("IMPOSSIBLE");
}
\end{code}

\usetemplate{Find a directed cycle}{Ctrl+Alt+T, /}{cpcycle}

% ======================================================================================================================
\section[Functional graphs]{Functional graphs \hfill \gold}\label{sec:functional}

In a \defn{functional graph} (or successor graph) every vertex has exactly one outgoing edge, $v \to f(v)$. They arise from
permutations, from ``each cow points to its best friend'' and from any deterministic process, where $f$ maps a state to the
next one (Section~\ref{sec:simulation}). Their structure is rigid. Follow the edges from any vertex: since there are finitely
many vertices, the walk eventually repeats a vertex and from then on goes around a cycle forever. So every component consists
of exactly one cycle with trees hanging into it, their edges pointing towards the cycle (Figure~\ref{fig:functional}).

\begin{figure}[htbp]
\centering
\begin{tikzpicture}
  \foreach \kk in {0,...,5} \node[vtx, cool] (c\kk) at ({2.2 + 1.15*cos(60*\kk)}, {1.15*sin(60*\kk)}) {};
  \foreach \kk in {0,...,5} {
    \pgfmathtruncatemacro{\nk}{mod(\kk + 1, 6)}
    \draw[->, line width=1.1pt, draw=fgblue] (c\kk) to[bend right=18] (c\nk);
  }
  \node[vtx] (t1) at (-0.35, 0) {}; \node[vtx] (t2) at (-1.55, 0.2) {}; \node[vtx] (t3) at (-1.2, 1.2) {};
  \node[vtx] (t6) at (-1.3, -1.0) {};
  \node[vtx] (t4) at (4.6, 1.35) {}; \node[vtx] (t5) at (5.7, 2.0) {};
  \draw[ded] (t2) -- (t1); \draw[ded] (t3) -- (t1); \draw[ded] (t6) -- (t1); \draw[ded] (t1) -- (c3);
  \draw[ded] (t5) -- (t4); \draw[ded] (t4) -- (c1);
  \node[font=\sffamily\scriptsize, text=fgblue] at (2.2, 0) {cycle};
  \node[notel, anchor=west, text width=5.2cm] at (6.3, 0.2) {Blue: the cycle of the component. Grey: trees hanging into
    it, all edges pointing towards the cycle. From any vertex, the walk reaches the cycle after some number of steps (the
    tail) and then repeats with the cycle's length.};
\end{tikzpicture}
\caption{A component of a functional graph: exactly one cycle, with in-trees attached to its vertices.}
\label{fig:functional}
\end{figure}

For every vertex, two numbers describe its future completely: the length of its tail (steps until it reaches a cycle) and the
length of that cycle. Both can be computed for all vertices in $\Oh{n}$ by walking from each unvisited vertex until the walk
hits a vertex that is either on the current walk (a new cycle has been found) or already solved (then the tail lengths follow
by counting back); Planets Cycles in Chapter~\ref{ch:worked} does exactly this. ``Where is $v$ after $k$ steps?'' for many
$v$ and huge $k$ is answered by binary lifting in $\Oh{\log k}$ each.

\usetemplate{Functional graph (successor graph) cycles}{Ctrl+Alt+T, ,}{cpfuncgraph}

\usetemplate{K-th successor (binary lifting / jumping)}{Ctrl+Alt+T, .}{cpjump}

\subsection{Cycle finding with constant memory}

Sometimes $f$ is given by a formula on a huge domain ($f(x) = x^2 + c \bmod N$ in Pollard's rho, the next state of a large
simulation) and the visited vertices cannot be stored. \defn{Floyd's cycle-finding algorithm} runs two walkers from $x_0$, the
tortoise one step at a time and the hare two steps at a time. Once both are on the cycle, the hare gains one step per round on
the tortoise, so they meet within $\mu + \lambda$ rounds, where $\mu$ is the tail length and $\lambda$ the cycle length.
\defn{Brent's algorithm} is the practical refinement: the hare walks alone while the tortoise waits at the positions
$1, 2, 4, 8, \ldots$; when the hare meets the waiting tortoise, the number of steps since the last teleport is exactly
$\lambda$. Then two walkers started $\lambda$ steps apart from $x_0$ meet precisely at the cycle's entry, after $\mu$ steps.
Brent uses about a third fewer evaluations of $f$ than Floyd, and both use $\Oh{1}$ memory.

\begin{code}
template <class T, class F>
pair<ll, ll> find_rho(T x0, F f) {  // returns (mu, lambda)
    ll power = 1, lam = 1;
    T tortoise = x0, hare = f(x0);
    while (!(tortoise == hare)) {  // the tortoise waits at 1, 2, 4, ...
        if (power == lam) tortoise = hare, power *= 2, lam = 0;
        hare = f(hare), lam++;
    }
    T a = x0, b = x0;  // b starts lambda steps ahead: they meet at the entry
    for (ll i = 0; i < lam; i++) b = f(b);
    ll mu = 0;
    while (!(a == b)) a = f(a), b = f(b), mu++;
    return {mu, lam};
}
\end{code}

\usetemplate{Cycle finding in O(1) memory (Floyd / Brent)}{Ctrl+Alt+U, F}{cpcyclefind}

% ======================================================================================================================
\section[Bidirectional search and A*]{Bidirectional search and A* \hfill \ioi}\label{sec:bibfs}

In an implicit graph with branching factor $b$, a BFS to depth $d$ explores about $b^d$ states. Two ideas cut this down, one
by halving the exponent and one by steering the search.

\paragraph{Bidirectional BFS.} When the target is known, search from both ends at once, expanding one whole layer at a time
from whichever side has the smaller frontier, and stop as soon as the two searches touch. Each side then only needs depth about
$d/2$, so the work drops from about $b^d$ to about $2b^{d/2}$ (Figure~\ref{fig:bibfs}). This is meet in the middle
(Section~\ref{sec:mitm}) for shortest paths. The first contact found while expanding a layer gives a shortest path: before that
layer, no contact means the distance exceeds the sum of the two explored depths, and every contact found during the layer has
exactly that length plus one.

\begin{figure}[htbp]
\centering
\begin{tikzpicture}
  \begin{scope}
    \fill[fgsky!25] (0, 0) circle (2.0);
    \draw[line width=0.8pt, draw=fgblue] (0, 0) circle (2.0);
    \fill[fgblue] (0, 0) circle (2.5pt) node[below, font=\sffamily\small] {$s$};
    \fill[fgverm] (2.0, 0) circle (2.5pt) node[right, font=\sffamily\small] {$t$};
    \node[lbl] at (0, -2.4) {one search: radius $d$};
  \end{scope}
  \begin{scope}[xshift=6.4cm]
    \fill[fgsky!25] (-1.0, 0) circle (1.0);
    \draw[line width=0.8pt, draw=fgblue] (-1.0, 0) circle (1.0);
    \fill[fgorange!25] (1.0, 0) circle (1.0);
    \draw[line width=0.8pt, draw=fgorange!85!black] (1.0, 0) circle (1.0);
    \fill[fgblue] (-1.0, 0) circle (2.5pt) node[below, font=\sffamily\small] {$s$};
    \fill[fgverm] (1.0, 0) circle (2.5pt) node[below, font=\sffamily\small] {$t$};
    \node[lbl] at (0, -2.4) {two searches: radius $d/2$ each};
  \end{scope}
\end{tikzpicture}
\caption{Why bidirectional search wins: in a graph where the number of states within distance $r$ grows like $b^r$, two balls
of radius $d/2$ contain far fewer states than one ball of radius $d$.}
\label{fig:bibfs}
\end{figure}

\begin{problem}{Swap Game (CSES)}
A $3 \times 3$ grid holds the numbers $1, \ldots, 9$. One move swaps two adjacent numbers. Find the minimum number of moves
that sorts the grid.
\solution The states are the $9! = 362\,880$ permutations and each has 12 neighbours, so even a plain BFS fits. Bidirectional
BFS, from the start and from the sorted grid, explores far fewer states whenever the answer is not close to the maximum
(Figure~\ref{fig:swapstates}). Each state is encoded as the nine digits read row by row, a number below $10^9$.
\end{problem}

\begin{code}
void solve() {
    // a state is the 3 x 3 grid read row by row as a 9-digit number
    auto encode = [](const array<int, 9>& a) {
        int x = 0;
        for (int v : a) x = x * 10 + v;
        return x;
    };
    auto decode = [](int x) {
        array<int, 9> a;
        per(i, 9) a[i] = x % 10, x /= 10;
        return a;
    };
    array<int, 9> start;
    for (int& v : start) cin >> v;
    int s = encode(start), t = 123456789;
    vector<pii> swaps;  // the 12 pairs of adjacent cells
    rep(i, 9) {
        if (i % 3 < 2) swaps.push_back({i, i + 1});
        if (i < 6) swaps.push_back({i, i + 3});
    }
    // bidirectional BFS: grow the smaller frontier until they touch
    unordered_map<int, int> dist[2];
    vi frontier[2] = {{s}, {t}};
    dist[0][s] = 0, dist[1][t] = 0;
    if (s == t) {
        print(0);
        return;
    }
    while (true) {
        int side = sz(frontier[0]) <= sz(frontier[1]) ? 0 : 1;
        vi next;
        for (int x : frontier[side]) {
            array<int, 9> a = decode(x);
            int d = dist[side][x];
            for (auto [i, j] : swaps) {
                swap(a[i], a[j]);
                int y = encode(a);
                swap(a[i], a[j]);
                if (dist[side].count(y)) continue;
                auto it = dist[side ^ 1].find(y);
                if (it != dist[side ^ 1].end()) {  // the two searches meet
                    print(d + 1 + it->second);
                    return;
                }
                dist[side][y] = d + 1;
                next.push_back(y);
            }
        }
        frontier[side] = next;
    }
}
\end{code}

\begin{figure}[htbp]
\centering
\begin{tikzpicture}
  \begin{axis}[cpplot, width=10.5cm, height=5.6cm, ymode=log, xmin=1, xmax=17, ymin=10, ymax=1e6,
      xlabel={distance $d$ between start and target}, ylabel={states explored}, xtick={2,4,...,16},
      legend pos=south east]
    \addplot coordinates {(2,101) (4,2549) (6,27288) (8,131470) (10,291158) (12,358449) (14,362855) (16,362880)};
    \addlegendentry{one BFS: states within $d$}
    \addplot coordinates {(2,26) (4,202) (6,1142) (8,5098) (10,18414) (12,54576) (14,132448) (16,262940)};
    \addlegendentry{bidirectional: twice the states within $d/2$}
  \end{axis}
\end{tikzpicture}
\caption{Swap Game: the number of states that each search explores in the worst case, computed from the exact distance
distribution of all $9!$ states. At distance 10 bidirectional search explores 16 times fewer states; near the maximum distance
16 the gain shrinks, because a plain BFS cannot explore more than the $9!$ states that exist.}
\label{fig:swapstates}
\end{figure}

\usetemplate{Bidirectional BFS (meet in the middle on shortest paths)}{Ctrl+Alt+U, 2}{cpbibfs}

\paragraph{A*.} When a good estimate of the remaining distance is available, a search can be \emph{steered} towards the target.
\defn{A*} is Dijkstra's algorithm (Chapter~\ref{ch:shortest}) with priority $g(v) + h(v)$, where $g(v)$ is the distance found so
far and $h(v)$ is a lower bound on the remaining distance to the target, such as the Manhattan distance on a grid. If $h$ never
overestimates (it is \emph{admissible}), the first time the target is taken from the priority queue its distance is optimal. If
moreover $h(u) \le w(u, v) + h(v)$ for every edge (it is \emph{consistent}, as the Manhattan distance is on a grid), no vertex
is ever expanded twice, and A* is exactly Dijkstra on the edge weights $w(u, v) - h(u) + h(v) \ge 0$. With $h = 0$ it is plain
Dijkstra; the better the bound, the fewer vertices are expanded (Figure~\ref{fig:astar}).

\begin{figure}[htbp]
\centering
\begin{tikzpicture}[x=0.36cm, y=0.36cm]
  \begin{scope}[xshift=0cm]
    \foreach \rw/\ci in {
      6/2, 5/2, 7/2, 6/1, 6/3, 4/2, 5/1, 5/3, 8/2, 7/1, 7/3, 6/0, 6/4, 3/2, 4/1, 4/3, 5/0, 5/4, 8/1, 8/3, 7/0, 7/4,
      6/5, 2/2, 3/1, 4/0, 4/4, 5/5, 8/0, 8/4, 7/5, 6/6, 1/2, 2/1, 2/3, 3/0, 4/5, 5/6, 8/5, 7/6, 6/7, 0/2, 1/1, 1/3,
      2/0, 2/4, 4/6, 5/7, 8/6, 7/7, 0/1, 0/3, 1/0, 1/4, 2/5, 4/7, 8/7, 0/0, 0/4, 1/5, 2/6, 0/5, 1/6, 2/7, 0/6, 1/7,
      0/7, 0/8, 0/9, 1/9, 0/10, 2/9, 1/10, 0/11, 3/9, 2/10, 1/11, 0/12, 4/9, 3/10, 2/11, 1/12, 0/13, 5/9, 4/10, 3/11,
      2/12, 1/13, 0/14, 6/9, 5/10, 4/11, 3/12, 2/13, 1/14, 0/15, 7/9, 6/10, 5/11, 4/12, 3/13, 2/14}
      \fill[fgsky!40] (\ci, -\rw) rectangle ++(1, -1);
    \foreach \rw/\ci in {1/8, 2/8, 3/3, 3/4, 3/5, 3/6, 3/7, 3/8, 4/8, 5/8, 6/8, 7/8, 8/8}
      \fill[fgink!75] (\ci, -\rw) rectangle ++(1, -1);
    \draw[fgline!60, line width=0.25pt, step=1] (0, 0) grid (17, -9);
    \draw[line width=1.3pt, draw=fgverm] (2.5, -6.5) -- (2.5, -0.5) -- (9.5, -0.5) -- (9.5, -2.5) -- (14.5, -2.5);
    \fill[fggreen!60!black] (2.5, -6.5) circle (4pt);
    \fill[fgverm] (14.5, -2.5) circle (4pt);
    \node[lbl] at (8.5, -9.9) {BFS: 102 cells expanded};
  \end{scope}
  \begin{scope}[xshift=6.84cm]
    \foreach \rw/\ci in {
      6/2, 5/2, 4/2, 3/2, 2/2, 2/3, 2/4, 2/5, 2/6, 2/7, 4/3, 4/4, 4/5, 4/6, 4/7, 5/3, 5/4, 5/5, 5/6, 5/7, 6/3, 6/4,
      6/5, 6/6, 6/7, 1/7, 1/6, 1/5, 1/4, 1/3, 7/7, 1/2, 2/1, 7/6, 3/1, 7/5, 4/1, 7/4, 5/1, 7/3, 6/1, 7/2, 0/7, 0/8,
      0/9, 0/10, 0/11, 0/12, 0/13, 0/14, 1/14, 2/14}
      \fill[fgsky!40] (\ci, -\rw) rectangle ++(1, -1);
    \foreach \rw/\ci in {1/8, 2/8, 3/3, 3/4, 3/5, 3/6, 3/7, 3/8, 4/8, 5/8, 6/8, 7/8, 8/8}
      \fill[fgink!75] (\ci, -\rw) rectangle ++(1, -1);
    \draw[fgline!60, line width=0.25pt, step=1] (0, 0) grid (17, -9);
    \draw[line width=1.3pt, draw=fgverm] (2.5, -6.5) -- (2.5, -2.5) -- (7.5, -2.5) -- (7.5, -0.5) -- (14.5, -0.5) -- (14.5, -2.5);
    \fill[fggreen!60!black] (2.5, -6.5) circle (4pt);
    \fill[fgverm] (14.5, -2.5) circle (4pt);
    \node[lbl] at (8.5, -9.9) {A* with the Manhattan distance: 52 cells};
  \end{scope}
\end{tikzpicture}
\caption{BFS versus A* from the green start to the red target around a wall (dark). Blue cells were expanded before the target;
both searches find a shortest path of length 20 (red), but A* with the Manhattan distance as its bound expands half as many
cells.}
\label{fig:astar}
\end{figure}

A* is the method of choice for puzzles with a good lower bound (the sum of Manhattan distances of the tiles in the 15-puzzle),
and its memory-light variant IDA*, iterative deepening on the bound $g + h$, solves the 15-puzzle within the limits of an
ordinary contest problem. In typical contest graphs without such a bound, bidirectional BFS or a smarter state space are
usually the better investment.

\begin{summary}[Chapter summary]
\begin{itemize}
  \item Store graphs as adjacency lists by default; matrices for small dense graphs, edge lists for edge-by-edge algorithms.
  \item DFS runs in $\Oh{n + m}$; its tree and discovery times classify edges into tree, back, forward and cross edges. In
        undirected graphs only tree and back edges exist.
  \item BFS finds shortest paths in unweighted graphs; several sources go into the queue at once; directions or parents
        recover the path.
  \item A graph is bipartite if and only if it has no odd cycle; BFS parity finds the colouring or the cycle.
  \item Undirected cycles are back edges to non-parent ancestors; directed cycles are edges into grey vertices.
  \item Every component of a functional graph is one cycle with in-trees; Floyd and Brent find tails and cycles in
        $\Oh{1}$ memory.
  \item Bidirectional BFS turns $b^d$ into $2b^{d/2}$; A* steers Dijkstra with an admissible lower bound.
\end{itemize}
\end{summary}

\begin{practice}
\ppgroup{Traversals}
\begin{pplist}
  \item Counting Rooms
  \item Labyrinth
  \item Building Roads
  \item Message Route
  \item Building Teams
  \item Monsters
\end{pplist}
\ppgroup{Cycles, functional graphs and searches}
\begin{pplist}
  \item Round Trip
  \item Round Trip II
  \item Planets Cycles
  \item Planets Queries I
  \item Swap Game
  \item Planets Queries II
\end{pplist}
\end{practice}
